\scanpage{12}

\begin{center}
U N I V E R S I T E\\
DES SCIENCES ET TECHNIQUES\\
DU LANGUEDOC\\
\textasciicircum\textasciicircum\textasciicircum\textasciicircum\textasciicircum\textasciicircum\textasciicircum\textasciicircum\textasciicircum\textasciicircum\textasciicircum\textasciicircum\textasciicircum
\end{center}

Montpellier, le 14 octobre 1983

Service de la Scolarité

Diffusion : Liste A\\
MA/gs \hfill N° 4470

\begin{center}
AVIS DE SOUTENANCE D'UNE THESE DE DOCTEUR-INGENIEUR

\uuline{`` PHYSIQUE DES MATERIAUX ''}
\end{center}

\uline{NOM et Prénom du candidat} : GOUDEAU Jacques

\uline{Titre de la thèse} : ``Contraintes de trempe thermique dans des pièces à symétrie de révolution''

\uline{Date, heure et lieu de soutenance prévus} : Le 28 octobre 1983, à 11 heures\\
Salle de l'IREM - Université des Sciences et Techniques du Languedoc

\uline{Nom du Directeur de thèse} : M. ZARZYCKI - M. MAISONNEUVE

\uline{Désignation du Laboratoire où la thèse est préparée} : Laboratoire de Sciences des Matériaux, Mécanique des Milieux continus, Chimie structurale - U.S.T.L.

\uline{Analyse de l'ouvrage} : On étudie l'évolution des champs de température et de contraintes, dans une pièce en verre à symétrie de révolution.

Pour cela on résout par éléments finis un problème de thermovisco-élasticité, avec une loi de comportement de type MAXWELL à coefficients variables en fonction de la température.

On compare les résultats obtenus à ceux d'une méthode optique d'analyse des contraintes ainsi qu'à une méthode industrielle de mésure.

\uline{Durée approximative des recherches ayant abouti à la thèse} : 2 ans

\uline{Composition du jury proposé} : \begin{tabular}[t]{ll}
MM. & ZARZYCKI J.\\
& MAISONNEUVE O.\\
& REGNIER J\\
& FRANCOIS\\
& BOURRET D.
\end{tabular}

\begin{flushright}
P/ le Secrétaire Général,\\
l'Attachée Principale d'Administration\\
Scolaire et Universitaire,

\unsure{A. Gentil}

\uline{A. GENTIL}
\end{flushright}

\margin{11}
